% Notes of the Montreal seminar, 1970 (cote 129, typed part): archive cote 129, site page 72 (PDF page 73).
% typescript. Draft compared: none.
% AI-transcribed by Claude Opus 5.5 (claude-opus-5-5, adaptive thinking, effort high), Message Batches API, batch mtl129_full (msgbatch_01AGnDKHwy6ATCM7ehUa6263), 2026-11-13; not yet checked by eye.
\scanpage{129/72}
\mspage{\uline{III}-17}

ces deux conditions sont équivalentes et équivalent à chacune des con-%
ditions (ne faisant intervenir que la fibre géométrique de $G$):

3) $G_{\bar{s}}$ est p-divisible.

4) Il est impossible de plonger le groupe formel additif $\overline{\mathbb{G}}_a$
(complété formel de $\mathbb{G}_a$ le long de sa section unité) dans $G_{\bar{s}}$ .

Dans le cas général, on pense que les conditions 1) et 2) sont
encore équivalentes; cela résulterait de la conjecture suivante, vraie
dans le cas artinien local (SGA 3, VII, \quad ):

Si $u\colon G \to G'$ est un S-morphisme de variétés de Lie formelles
de même dimension relative sur $S$, on a les équivalences:

- $u$ est un épimorphisme

- Ker $(u)$ est représentable par un schéma fini

- Ker $(u)$ est représentable par un schéma fini localement libre.

Par contre, le fait que le groupe de Lie formel $G$ soit un
groupe de BT ne se vérifie pas en général sur les fibres, comme le mon-%
tre l'exemple du groupe de Lie formel associé à la courbe elliptique mo-%
dulaire $A$ sur une courbe $S$ de caractéristique $p$, puisqu'il y a des
points $s \in S$ tels que l'invariant de Hasse de $A_s$ est nul.

Cet exemple suggère cependant la conjecture suivante: Pour
qu'un groupe de Lie formel $G$ sur $S$ soit de BT, il faut et il suffit
que ses fibres $G_s$ soient des groupes de BT et que le rang de $G_s$
(en tant que BT) soit une fonction localement constante en $s \in S$.

\mspage{\hand{72}}
